%% Os labels são definidos por quem está escrevendo
\section{Specific requirements}
\subsection{Funcionais}    
    \subsubsection{Certificado} 
        \subsubsubsection{O sistema deve emitir certificados com o nome do participante, o tipo de evento que participou com o nome, ministrante ou palestrante, data, hora, e carga horária. Caso seja evento grande, será listado apenas aqueles que o usuário participou, com as datas, horário, carga horária e uma autenticação digital.}\label{c1}
        \subsubsubsection{O sistema deve emitir certificados para ministrantes ou palestrantes após a realização do evento com as seguintes informações: nome do evento, data, hora, carga horária, tipo de evento, local, nome do coordenador do evento e uma autenticação digital.}\label{c2}
        \subsubsubsection{O sistema deve emitir certificados, ou listar a participação de em um curso dentro de um grande evento, apenas para aqueles que participaram integralmente de todas as partes de cursos (caso seja mais que uma parte).}\label{c3}
        \subsubsubsection{Deve registrar e armazenar o nome e o diretório que está armazenado para a reemissão futura, também o registro da data e horário da emissão.}\label{c4}
        \subsubsubsection{O sistema deve entregar o certificado em formato de PDF.}\label{c5}
        % \subsubsubsection{O sistema deve emitir certificados somente para os usuários que participaram de todas as partes do curso ou palestra}
        
    \subsubsection{Usuários}
        \subsubsubsection{O sistema deve cadastrar usuários que não possuem cadastro com essas informações: nome, e-mail, tipo de vinculo com a universidade.}\label{u1}
        \subsubsubsection{A autenticação deve ser utilizando o serviço da \textit{Google}}\label{u2}
        \subsubsubsection{Caso o usuário seja externo à faculdade e não possua e-mail institucional, deve-se permitir o cadastro com um e-mail pessoal do usuário.}\label{u3}
        \subsubsubsection{O sistema deve permitir um usuário se tornar um palestrante ou ministrante de curso mediante a autorização do moderador do sistema.}\label{u4}
        \subsubsubsection{O login do usuário deve persistir no dispositivo em que entrou, não sendo necessário a sua autenticação constante.}\label{u5}
        \subsubsubsection{O sistema deve armazenar em quais eventos os usuários estão inscritos.}\label{u6}
        \subsubsubsection{O sistema deve permitir um usuário se tornar um moderador ou coordenador de eventos mediante a concessão de um coordenador de evento.}\label{u6}
    
    \subsubsection{Grande evento}
        \subsubsubsection{O sistema deve armazenar o nome, tema, cursos e palestras que faz parte do grande evento.}\label{g1}
        \subsubsubsection{Um grande evento deve ser criado somente por coordenadores de evento, moderadores do sistema, ministrantes ou palestrantes.}\label{g2}
        \subsubsubsection{O sistema deve permitir cadastro de eventos grandes com taxa de inscrição.}\label{g3}
        \subsubsubsection{Caso o grande evento possua limitação de vagas disponíveis, é necessário que o sistema manipule a quantidade de vagas disponíveis.}\label{g4}
        \subsubsubsection{Um grande evento é um agrupamento de cursos e palestras que os usuários participantes podem se inscrever, pode haver limitação de inscrições em palestras e cursos por usuário.}\label{g5}
    
    \subsubsection{Cursos}
        \subsubsubsection{O curso precisa armazenar o ministrante(s), título, tema, data, hora, local, requisito de instalação de softwares nos computadores do local do curso, quantidade de vagas disponíveis e quantidade de partes que o curso possui.}\label{cu1}
        \subsubsubsection{O ministrante é um usuário do sistema que possui cargo de ministrante dentro do sistema.}\label{cu2}
        \subsubsubsection{O curso deve ser criado somente por ministrantes, coordenador de evento ou moderador do sistema.}\label{cu3}
        \subsubsubsection{O sistema deve permitir o cadastro de curso com taxa de inscrição.}\label{cu4}
        \subsubsubsection{Caso o curso possua limitação de vagas disponíveis, é necessário que o sistema manipule a quantidade de vagas disponíveis.}\label{cu5}
    
    \subsubsection{Palestras}
        \subsubsubsection{A palestra precisa armazenar o palestrante, data, hora, local, capacidade máxima, tema e título.}\label{p1}
        \subsubsubsection{O palestrante é um usuário do sistema que possui cargo de ministrante dentro do sistema.}\label{p2}
        \subsubsubsection{A palestra deve ser criado somente por palestrantes, coordenador de evento ou moderador do sistema.}\label{p3}
        \subsubsubsection{O sistema deve permitir o cadastro de palestras com taxa de inscrição.}\label{p4}
        \subsubsubsection{Caso a palestra possua limitação de vagas disponíveis, é necessário que o sistema manipule a quantidade de vagas disponíveis.}\label{p5}
        
    \subsubsection{Inscrições em cursos, palestras e grandes eventos}
        \subsubsubsection{O sistema deve realizar o cadastro de inscrições de usuários em eventos grandes, cursos ou palestras.}\label{i1}
        % \subsubsubsection{O sistema deve permitir os participantes a modificação de suas respectivas inscrições em cursos ou palestras.}\label{i2}
        \subsubsubsection{O sistema deve armazenar a data e hora da inscrição dos usuários em cursos, palestras e grandes eventos.}\label{i3}
        \subsubsubsection{O sistema não deve permitir a inscrição em um curso ou palestra de um usuário caso possua conflito entre outros cursos ou palestras, seja em grande evento ou não.}\label{i4}
        \subsubsubsection{O sistema não deve permitir o pagamento da inscrição em palestras, cursos ou grandes eventos após 24h}\label{i5}
        \subsubsubsection{O sistema deve permitir a modificação de inscrição de palestra e ou cursos da sua respectiva inscrição em grande evento.}\label{i6}
        \subsubsubsection{O sistema deve permitir que um usuário cancele a sua inscrição em um curso ou palestra, assim liberando a vaga que estava ocupando.}\label{i7}    
        \subsubsubsection{O sistema deve permitir que um usuário cancela a sua inscrição em um grande evento, assim, liberando a vaga que estava ocupando.}\label{i8}
    
    \subsubsection{Presença}
        \subsubsubsection{O sistema deve armazenar presenças de palestras e/ou cursos, as partes e os participantes presentes durante a sua realização.}\label{pr1}
        \subsubsubsection{O sistema deve armazenar presenças de todos os eventos (incluindo todas as partes que pertencem ao evento) presentes em um grande evento.}\label{pr2}
        
    \subsubsection{Pagamento de taxa de inscrição em cursos, palestras e grandes eventos}
        \subsubsubsection{O sistema deve emitir o \textit{QR code} de pagamento por \textit{PIX} para cursos, palestras ou eventos grandes que são pagos, aprovar a inscrição automaticamente após a confirmação do pagamento ou excluir a inscrição caso o participante não pague dentro do período estipulado.}\label{pa1}
        \subsubsubsection{Caso o usuário cancele a sua inscrição, se a palestra, curso ou grande evento possua taxa de inscrição, a taxa deve ser estornado ao usuário.}\label{pa2}        

\subsection{Não funcionais}
    \subsubsection{Relatórios}
        \subsubsubsection{O sistema deve emitir um relatório, usando como filtro um intervalo entre as datas, buscando todos os cursos e palestras que ocorreram durante esse período, listando os nomes, datas, carga horária e ministrante(s) ou palestrante(s).}\label{r1}
\newpage